% ---
% id: meat-core/frank-wolfe/rem:fo-iterate-repair
% kind: rem
% title: "A repaired iterate as a construction step"
% summary: >-
%   A block repaired while the iterate is formed makes the repaired point the
%   iterate, at which every term of the bracket is evaluated, so no
%   perturbation that (H-NOPERT) forbids occurs. The repair moves the coupling
%   residual by at most the dimension times its size and says nothing about
%   the backend.
% prominence: standard
% status: accepted
% uses: [thm:firstorder, rem:fo-nopert, rem:fo-arithmetic]
% mentions: [req:iterate-checks, req:iterate-repair, foundations/req:norm-conversion, rem:fo-iterate-checks, rem:fo-oracle-repair, rem:fo-check-exposures]
% assumes: []
% refs: []
% provenance: newly-derived
% lean:
% ---

\begin{remark}\label{rem:fo-iterate-repair}
\emph{On a formed iterate: a construction step, and the only repair this
module licenses on an iterate.} The point consumed at iterate $t$ is formed
by the engine --- in the intended use a convex combination of oracle
returns --- and Requirement~\ref{req:iterate-checks} requires it to be
per-block positive semidefinite at exactly $0$; a block failing that check
is rejected or projected (Requirement~\ref{req:iterate-repair}).
\emph{The repaired point is $x_t$.} The repair belongs to the formation of
the iterate, upstream of the point the theorem is applied at, and it is a
different act from the oracle repair of Remark~\ref{rem:fo-oracle-repair},
which touches only the search; everything in \eqref{eq:hlb} --- the
objective $O^{J}_{\alpha}(x_t)$, the gradient $(\lmoq_t,\lmoJ_t)$ and the
pairings with $x_t$ --- is then evaluated at that one point. \emph{This is not the perturbation route
Remark~\ref{rem:fo-nopert} closes, and (H-NOPERT) is not violated}: that
route takes the gradient at a perturbed operator and concludes something
about the objective at the unperturbed one, which needs a domination
statement this module does not prove, whereas here no second function
enters the argument at all --- Theorem~\ref{thm:firstorder} concludes about
the infimum over the whole set and asks of $x_t$ only that it lie in
$\acSet\times\Feas^{J}(\varepsilon)$. What (H-NOPERT) forbids is a
perturbation applied \emph{to} $x_t$, between the point being fixed and the
gradient being taken; a repair applied \emph{before} the point is fixed is a
different act, and it is the one Requirement~\ref{req:iterate-repair}
licenses.

\emph{The repair moves the residuals by a bounded amount.} A projection
moves the coupling and $\bowtie_k$ residuals of
Remark~\ref{rem:fo-iterate-checks}: if the repair moves each block by at
most $\delta$ in the operator norm, the coupling residual moves by at most
$d_{\Qsys}\,\delta$ in $\|\cdot\|_{\max}$, since
$\|\tr_{\Qsys_s}\Delta_s\|\le d_s\|\Delta_s\|$, $\sum_sd_s=d_{\Qsys}$ and
$\|\cdot\|_{\max}\le\|\cdot\|$, a bound that carries its dimension factor
as \xref{foundations/req:norm-conversion} asks. The requirement nevertheless
measures the residuals again rather than bounding them in advance.
What the repair establishes is a hypothesis of
Theorem~\ref{thm:firstorder}, on an object the engine itself chose, and it
says nothing whatever about the backend.

\emph{What it does not establish.} $\lmin$ is itself computed, so a block
whose true minimum eigenvalue is negative within the eigensolver's own
backward error passes the check unrepaired, and the reassembly of a
projected block is itself rounded. That residue is a third one-sided
exposure of the kind Remark~\ref{rem:fo-check-exposures} names, and it is
named there with the other two. Its \emph{nature} is the unbudgeted
arithmetic this module declares everywhere (Remark~\ref{rem:fo-arithmetic}),
and not a window this module chose while knowing it admitted points at which
its own theorem does not apply.
\end{remark}
